% Generated draft: requires mathematical and source review.
\documentclass[11pt]{article}
\usepackage[margin=25mm]{geometry}
\usepackage{fontspec}
\setmainfont{Latin Modern Roman}
\usepackage{amsmath,amssymb,mathtools,unicode-math}
\setmathfont{Latin Modern Math}
\usepackage{enumitem,needspace,xurl,hyperref}
\setlength{\parindent}{0pt}
\setlength{\parskip}{5pt}
\setlength{\emergencystretch}{3em}
\newcommand{\N}{\mathbb{N}}
\newcommand{\Z}{\mathbb{Z}}
\newcommand{\Q}{\mathbb{Q}}
\newcommand{\R}{\mathbb{R}}
\newcommand{\C}{\mathbb{C}}
\newcommand{\F}{\mathbb{F}}
\newcommand{\A}{\mathbb{A}}
\newcommand{\PP}{\mathbb{P}}
\newcommand{\E}{\mathbb{E}}
\newcommand{\eps}{\varepsilon}
\newcommand{\dd}{\,\mathrm{d}}
\begin{document}
\section*{UnsolvedMath ID: 20002840}
\textbf{Source problem code:} catalogue entry
\subsection{Definitions and mathematical statement}
The source is Conjecture 5 in the section ``Taylor: Florian Herzig's Thesis'' of the February 2006 workshop notes. The supplied verification restores \(p>2\) and arithmetic Frobenius. Let \(k=\overline{\F}_p\), with any finite-field descent specified separately, and let \(\rho:G_\Q\to\mathrm{GL}_2(k)\) be a continuous absolutely irreducible odd representation. These hypotheses are needed background, not fully restated in the extracted conjecture. Let \(N\ge1\), \(p\nmid N\), and let
\[
 N(\rho)=\prod_{\ell\ne p}\ell^{a_\ell(\rho)}
\]
be its prime-to-\(p\) Artin conductor. Write \(\operatorname{JH}(\sigma)\) for the isomorphism classes of Jordan--H\"older constituents of a finite-dimensional modular \(\mathrm{GL}_2(\F_p)\)-representation \(\sigma\), and \(\operatorname{Wt}(\rho|_{I_p})\) for the intended set of Serre weights. The source does not give that local weight recipe.

Using the source's \(\Gamma_1(N)\), \(T_\ell\), and coefficient-action conventions, the printed assertion is
\[
 H^1(\Gamma_1(N),\sigma)\text{ contains a simultaneous Hecke eigenclass with }
 T_\ell\text{-eigenvalue }\operatorname{tr}\rho(\operatorname{Frob}_\ell)
 \text{ for every prime }\ell\nmid Np
\]
if and only if
\[
 N(\rho)\mid N,\qquad
 \operatorname{JH}(\sigma)\cap\operatorname{Wt}(\rho|_{I_p})\ne\varnothing.
\]
Determinant/nebentype, dual versus direct coefficient actions, irreducibility/localization hypotheses and the extension from irreducible weights to arbitrary \(\sigma\) must be specified before claiming this literal formulation as a theorem. The aid reports the core \(\mathrm{GL}_2/\Q\) Serre assertion as established by Khare--Wintenberger with Kisin's completion, subject to those conventions; it does not independently verify every compressed coefficient-system statement above.

The source asks how to generalize, records the possible inclusion \(\operatorname{Wt}(\rho|_{I_p})\subset\operatorname{Wt}(\rho|_{I_p}^{\mathrm{ss}})\), and mentions an earlier explanation by ``Matt'', whose identity and explanation are absent. Full inertial extension data and semisimplification need not be interchangeable. For the Ash et al. higher-rank discussion, let \(\rho:G_\Q\to\mathrm{GL}_n(k)\) be irreducible with complex conjugation \(c\) satisfying
\[
 \left|\dim\ker(\rho(c)-1)-\dim\ker(\rho(c)+1)\right|\le1.
\]
The intended weight set consists of irreducible modular \(\mathrm{GL}_n(\F_p)\)-representations, not a single irreducible representation. The block-matrix congruence in the source is naturally interpreted as
\[
 \Gamma_1^{(n)}(N)=\{g\in\mathrm{SL}_n(\Z):
          \text{last row of }g\equiv(0,\ldots,0,1)\pmod N\}.
\]
The general weight recipe, cohomological degree, coefficient system and converse are not completely stated; the aid reports the broad higher-rank weight/level/degree problem unresolved.

\textbf{Local Hecke-information correction in the aid.} The source displays only the double coset of \(\operatorname{diag}(\ell,1,\ldots,1)\), denoted \(T(\ell,1)\). In the Ash--Doud--Pollack normalization, a full unramified packet \(a(\ell,k)\) is encoded by
\[
 \det(1-\rho(\operatorname{Frob}_\ell)X)
 =\sum_{k=0}^n(-1)^k\ell^{k(k-1)/2}a(\ell,k)X^k,
 \qquad a(\ell,0)=1.
\]
The standard operators \(T(\ell,k)\) use double cosets with \(k\) diagonal entries equal to \(\ell\) and the others 1, in the compatible convention. The scalar powers of \(\ell\) are invertible modulo \(p\). The first operator gives trace, and the last coefficient gives determinant.

For \(n\ge3\), one trace and a fixed determinant do not determine an arbitrary local unramified characteristic polynomial. The aid's characteristic-independent matrices are the companions
\[
 C_0=\begin{pmatrix}0&0&1\\1&0&0\\0&1&0\end{pmatrix},\qquad
 C_1=\begin{pmatrix}0&0&1\\1&0&-1\\0&1&0\end{pmatrix}.
\]
They both have trace zero and determinant one, but characteristic polynomials \(Y^3-1\) and \(Y^3+Y-1\). For larger \(n\), adjoining \(I_{n-3}\) preserves equality of traces and determinants and the difference of characteristic polynomials. Thus, within the standard elementary-symmetric Hecke family, local polynomial recovery requires all \(T(\ell,k)\), \(1\le k\le n\), or \(1\le k\le n-1\) if determinant is already fixed. Omitting another independent coefficient leaves a local ambiguity.

This is a per-prime information test, not a counterexample to a global conjecture imposing trace identities at every Frobenius or additional global constraints. It neither constructs global Galois representations realizing the displayed local matrices nor settles the higher-rank Serre-weight problem.
\subsection{Short English statement}
The record combines a compressed \(\mathrm{GL}_2\) Serre-weight statement with a higher-rank generalization. The aid reports the core rank-two theorem subject to missing conventions and identifies a local limitation of the displayed trace operator in higher rank. Recovering an arbitrary local characteristic polynomial needs the remaining elementary-symmetric Hecke data, while the global higher-rank weight and degree questions remain separate.
\subsection{Sources}
\begin{itemize}
\item Local UnsolvedMath registry, permanent ID 20002840, catalogue entry; canonical \texttt{aim-representation-theory-notes.json}, notes[111].
\item AIM \emph{\(p\)-adic representations, modularity, and beyond}, February 2006, section 12 ``Taylor: Florian Herzig's Thesis'', Conjecture 5 and its higher-rank discussion; \url{https://aimath.org/WWN/padicmodularity/padicmodularity.pdf}. The supplied verification restores \(p>2\) and arithmetic Frobenius; it is not independently repeated here.
\item Source attributions: Ash et al.; an unidentified earlier ``Matt'' explanation. Aid attributions: Khare--Wintenberger, Kisin; Ash--Doud--Pollack normalization. No full bibliographic entries or theorem numbers are supplied.
\item Machine-generated research attempt 1, partial-result/partially-solved labels, reported rigor proved and low novelty confidence. Its local information obstruction is not peer reviewed and requires independent verification.
\end{itemize}
\end{document}
